\chapter{Harness, Permissions, and Diagnostics Journey}
\label{ch:harness}

\section{Execution Progression}

The execution model progressed through three practical stages:
\begin{enumerate}
  \item repo-local execution,
  \item devcontainer execution,
  \item host-integrated execution for real runtime diagnostics.
\end{enumerate}

\section{Diagnostics Loop}

\begin{figure}[h]
\centering
\begin{tikzpicture}[
  node distance=1.5cm,
  box/.style={draw, rounded corners, align=center, inner sep=5pt},
  >=Latex
]
\node[box] (a) {Capture failure};
\node[box, right=of a] (b) {Collect traces};
\node[box, right=of b] (c) {Classify domain};
\node[box, right=of c] (d) {Patch + tests};
\node[box, right=of d] (e) {Update docs};
\draw[->] (a) -- (b);
\draw[->] (b) -- (c);
\draw[->] (c) -- (d);
\draw[->] (d) -- (e);
\end{tikzpicture}
\caption{Repeatable troubleshooting loop used across incident classes.}
\end{figure}

\section{Incident Families Addressed}

\begin{itemize}
  \item constrained-host tool discovery and naming mismatches,
  \item MCP-Apps runtime and bridge behaviors,
  \item map payload and startup pressure,
  \item NOMIS/ONS query-shape drift,
  \item collection aliasing and count semantics in feature queries.
\end{itemize}

The resulting evidence base is represented in troubleshooting reports and tracker
workstreams \cite{troubleshooting_master,peat_failure_chain,nomis_deep_analysis}.
